\section{Introduction}

Artificial intelligence (AI) is increasingly being used to assist and automate decision-making across a wide range of domains, from financial advising to medical diagnosis and hiring decisions. As AI systems become more sophisticated and ubiquitous, understanding when and why people choose to delegate decisions to algorithms has become a critical question for researchers, practitioners, and policymakers \citep{lee2004, kohn2021}.

One common assumption in the design of AI systems is that personalization increases trust and encourages users to delegate decisions to algorithms \citep{longoni2019, jussupow2020}. Personalized algorithms that incorporate individual preferences, characteristics, or past behavior are often assumed to be more trustworthy and more likely to be adopted than one-size-fits-all algorithms. This assumption has led to widespread investment in personalization across industries, from personalized recommendations on e-commerce platforms to personalized financial advice from robo-advisors.

However, this assumption has not been systematically tested in the context of delegation to algorithms. While personalization may increase trust in some contexts \citep{longoni2019}, it may also have unintended consequences. For example, personalized algorithms may be perceived as less objective or less competent than algorithms that use objective criteria \citep{castelo2019}. They may also raise concerns about privacy, surveillance, or manipulation \citep{hoff2015}. The net effect of personalization on delegation decisions is therefore an open empirical question.

This study addresses this gap by examining how personalization affects delegation to algorithms in financial decision-making. We conducted a preregistered online experiment in which participants made financial decisions with the option to delegate to either an objective expected value (EV)-maximizing algorithm or a personalized algorithm that incorporated their individual preferences. Counterintuitively, we found that participants initially delegated more to the objective EV algorithm than to the personalized algorithm, even though the personalized algorithm should theoretically lead to higher expected utility for them.

This finding challenges common assumptions about the benefits of personalization in AI systems and has important implications for the design of algorithmic advisors. It suggests that in certain contexts, particularly those involving objective calculations, people may prefer algorithms that are perceived as competent and objective rather than those that attempt to personalize recommendations.

\subsection{Theoretical Background}

\subsubsection{Algorithm Appreciation vs. Aversion}

The literature on algorithm appreciation and aversion provides a useful starting point for understanding delegation to algorithms. \citet{dietvorst2015} introduced the concept of algorithm aversion, showing that people avoid algorithms after observing them make errors, even when the algorithms outperform humans. In contrast, \citet{logg2019} found evidence of algorithm appreciation, showing that people prefer algorithmic judgment over human judgment in objective tasks.

These findings suggest that algorithm aversion/appreciation depends on task type and context. \citet{castelo2019} showed that people are more likely to use algorithms for objective tasks (e.g., forecasting) than for subjective tasks (e.g., hiring). This aligns with the MABA-HABA (Machines Are Better At vs. Humans Are Better At) framework proposed by \citet{glikson2020}, which suggests that people are more willing to delegate tasks that machines are perceived to be better at.

\subsubsection{Trust and Reliance in AI}

Trust is a key factor influencing delegation to algorithms \citep{lee2004, kohn2021}. \citet{hoff2015} identified several factors that influence trust in automation, including perceived competence, reliability, and understandability. \citet{kohn2021} emphasized the importance of distinguishing between trust as an attitude and reliance as behavior, noting that people may trust an algorithm but still choose not to rely on it.

Personalization may affect trust in several ways. On the one hand, personalization may increase trust by making algorithms more relevant and useful to individual users \citep{longoni2019}. On the other hand, personalization may decrease trust by making algorithms seem less objective or less competent \citep{castelo2019}. The net effect of personalization on trust is therefore an open question.

\subsubsection{Delegation in Decision-Making}

Delegation is a fundamental aspect of decision-making, both in human-human and human-AI interactions \citep{freer2023, ivanovastenzel2025}. People delegate decisions for various reasons, including lack of expertise, time constraints, or a desire to avoid responsibility \citep{freer2023}. In the context of AI, delegation decisions are influenced by factors such as perceived algorithm competence, task characteristics, and individual differences \citep{holzmeister2023, gazit2023}.

\citet{holzmeister2023} examined delegation to financial advisors and found that people are more likely to delegate when they perceive the advisor as competent and trustworthy. \citet{gazit2023} found that people are more likely to delegate to algorithms than to human advisors for objective tasks, but not for subjective tasks. These findings suggest that the type of task and the perceived competence of the advisor are key factors influencing delegation decisions.

\subsection{Research Gap}

Despite the growing literature on algorithm appreciation, trust in AI, and delegation, there is a gap in our understanding of how personalization affects delegation to algorithms. While personalization is often assumed to increase trust and adoption \citep{longoni2019, jussupow2020}, this assumption has not been systematically tested in the context of delegation.

Moreover, most studies on algorithm appreciation and aversion have focused on comparing algorithms to humans, rather than comparing different types of algorithms to each other \citep{dietvorst2015, logg2019}. This study fills this gap by comparing delegation to an objective EV-maximizing algorithm versus a personalized algorithm that incorporates individual preferences.

\subsection{Research Question and Hypotheses}

This study addresses the following research question: How does personalization affect delegation to algorithms in financial decision-making?

We preregistered the following hypothesis:

\textbf{H1}: Participants will delegate more to a personalized algorithm that incorporates their individual preferences than to an objective EV-maximizing algorithm.

This hypothesis is based on the assumption that personalization increases trust and makes algorithms more relevant and useful to individual users \citep{longoni2019, jussupow2020}. However, as we will show, the results challenge this assumption.

\subsection{Contributions}

This study makes several contributions to the literature:

First, it provides novel evidence that personalization can deter rather than encourage delegation to algorithms, challenging common assumptions about the benefits of personalization in AI systems.

Second, it contributes to the literature on algorithm appreciation and aversion by showing that the type of algorithm matters, not just whether an algorithm is used \citep{dietvorst2015, logg2019}.

Third, it contributes to the literature on trust in AI by examining how personalization affects trust and reliance \citep{lee2004, kohn2021}.

Fourth, it has practical implications for the design of algorithmic advisors, particularly in financial advising \citep{holzmeister2023}.

Finally, it is a preregistered experiment, which increases the credibility of the findings and contributes to the growing movement toward open and transparent research practices.

\subsection{Outline}

The remainder of this paper is organized as follows. Section~\ref{sec:literature} reviews the relevant literature on algorithm appreciation, trust in AI, and delegation. Section~\ref{sec:method} describes the experimental design and methods. Section~\ref{sec:results} presents the results. Section~\ref{sec:discussion} discusses the theoretical and practical implications, limitations, and directions for future research. Section~\ref{sec:conclusion} concludes.